% !TeX program = xelatex
% Compile from this directory: latexmk -xelatex template.tex
\documentclass[12pt,a4paper]{article}
\usepackage[margin=2.54cm]{geometry}
\usepackage{fontspec}
% Times New Roman matches the Word reference; the fallback works on Overleaf.
\IfFontExistsTF{Times New Roman}{\setmainfont{Times New Roman}}{\setmainfont{TeX Gyre Termes}}
\usepackage{amsmath}
\usepackage{graphicx}
\usepackage{booktabs,tabularx}
\usepackage{caption}
\usepackage{enumitem}
\usepackage{titlesec}
\usepackage{setspace}
\usepackage{fancyhdr}
\usepackage[hidelinks]{hyperref}
\graphicspath{{fig/}}
\setstretch{1.25}
\setlength{\parindent}{2em}
\setlength{\parskip}{4pt}
\widowpenalty=10000
\clubpenalty=10000
\titleformat{\section}{\bfseries\fontsize{14}{21}\selectfont}{\thesection}{0.6em}{}
\titleformat{\subsection}{\bfseries\fontsize{12}{18}\selectfont}{\thesubsection}{0.6em}{}
\titlespacing*{\section}{0pt}{14pt}{6pt}
\titlespacing*{\subsection}{0pt}{10pt}{4pt}
\captionsetup{font=small,labelsep=period}
\captionsetup[table]{position=top}
\setlist[enumerate]{label=(\arabic*),leftmargin=2em,itemsep=3pt,topsep=4pt}
\hypersetup{pdftitle={Laboratory Report Template},pdfauthor={},pdfsubject={Reusable laboratory report template}}

\setlength{\headheight}{26pt}
\setlength{\headsep}{16pt}

% ---------- Cover fields: replace the values below ----------
\newcommand{\InstitutionName}{South University of Science and Technology of China}
\newcommand{\DepartmentName}{Department of Materials science and Engineering}
\newcommand{\CourseTitle}{Comprehensive Experiments of Materials I}
\newcommand{\ExperimentTitle}{Lab [number]: [Experiment title]}
\newcommand{\ClassNumber}{[Class number]}
\newcommand{\ListNumber}{[List number]}
\newcommand{\StudentName}{[Your name]}
\newcommand{\StudentID}{[Your student ID]}
\newcommand{\TeamMembers}{[Names of team members]}
\newcommand{\ExperimentDate}{[YYYY.MM.DD]}
\newcommand{\SubmissionDate}{[YYYY.MM.DD]}

\begin{document}

% ==============================================================================
%                           封面部分开始 / COVER START
% ==============================================================================
% Editable reconstruction of the cover in Lab report template.docx.
\hypersetup{pageanchor=false}
\begingroup
\newgeometry{margin=2.54cm,headheight=26pt,headsep=16pt}
\setlength{\headwidth}{\textwidth}
\fancypagestyle{cover}{%
    \fancyhf{}
    \fancyhead[C]{\bfseries\fontsize{10}{12}\selectfont
        \InstitutionName\\\DepartmentName}
    \renewcommand{\headrulewidth}{0.4pt}
}
\begin{titlepage}
    \thispagestyle{cover}
    \linespread{1.25}\selectfont
    \setlength{\parskip}{4pt}
    \centering
    \vspace*{2.3cm}
    {\bfseries\fontsize{18}{24}\selectfont \CourseTitle\par}
    \vspace{1.7cm}
    {\bfseries\fontsize{16}{22}\selectfont \ExperimentTitle\par}
    \vspace{4.4cm}
    \begin{flushright}
    \begin{minipage}{0.60\textwidth}
        \raggedright
        \bfseries\fontsize{14}{21}\selectfont
        Class No.: \ClassNumber\par
        Name list number: \ListNumber\par
        \vspace{1.5cm}
        Name: \StudentName\par
        Student ID: \StudentID\par
        Team member: \TeamMembers\par
        Date of experiment: \ExperimentDate\par
        Date of submission: \SubmissionDate\par
    \end{minipage}
    \end{flushright}
    \vfill
\end{titlepage}
\restoregeometry
\endgroup
\setcounter{page}{1}
\hypersetup{pageanchor=true}

% ==============================================================================
%                           封面部分结束 / COVER END
% ==============================================================================

\section*{Abstract}
Write a concise summary of the experiment. State the objective, the main methods,
the principal quantitative results, and the conclusion. Replace this guidance
with a self-contained paragraph after completing the report.

\section{Introduction}
\subsection{Background and principles}
Explain the scientific background and the principles needed to understand the
experiment. Define important terms and symbols, and cite relevant sources using
numbered citations such as Ref.~\cite{reference-placeholder}. Introduce each
figure in the text before it appears, as illustrated in Figure~\ref{fig:background}.

\begin{figure}[htbp]
    \centering
    \includegraphics[width=0.65\linewidth]{figA.jpg}
    \caption{Replace with a description of the background diagram or principle.}
    \label{fig:background}
\end{figure}

\subsection{Objective}
State the experimental objective or question. Explain which variables will be
controlled or compared, what will be measured, and how the measurements address
the objective.

\section{Materials and Apparatus (or Equipment)}
\subsection{Materials}
List the materials, reagents, and samples used. Include relevant grades,
concentrations, quantities, and suppliers where needed for reproducibility.
Replace the entries below with the actual materials.
\begin{itemize}[leftmargin=2em,itemsep=0pt]
    \item {}[Material or reagent name], [grade or purity], [quantity or concentration].
    \item {}[Sample name], [composition or preparation details].
\end{itemize}


\subsection{Apparatus}
List the instruments and equipment. Give the manufacturer, model, measurement
range, or resolution where relevant. Replace the entries below with the actual
apparatus.
\begin{itemize}[leftmargin=2em,itemsep=0pt]
    \item {}[Instrument name], [manufacturer and model].
    \item {}[Equipment name], [relevant specification].
\end{itemize}

\section{Experimental Procedures}
\subsection{Sample preparation}
Describe the work in chronological order with enough detail to reproduce it.
Record quantities, conditions, durations, and any departures from the prescribed
method. Replace these example steps and the diagram in Figure~\ref{fig:procedure}.
\begin{enumerate}
    \item Describe the initial preparation, including material quantities and conditions.
    \item Describe the processing, treatment, and collection of the sample.
\end{enumerate}

\begin{figure}[htbp]
    \centering
    \includegraphics[width=0.65\linewidth]{figB.jpg}
    \caption{Replace with a description of the experimental workflow or setup.}
    \label{fig:procedure}
\end{figure}

\subsection{Characterization and measurement}
Describe sample handling, instrument settings, calibration, measurement
conditions, and the number of repeats. Explain how raw measurements are
converted into the reported quantities. Define every symbol and include units.
The equation below demonstrates a numbered equation and cross-reference:
\begin{equation}
    y = ax + b.
    \label{eq:example}
\end{equation}
Replace Equation~\eqref{eq:example} with the experimental relation. Here, $x$ and
$y$ are the input and response, $a$ is the slope, and $b$ is the intercept.
The equation is a formatting example only.

\section{Results and discussion}
\subsection{Observations and characterization}
Present the main observations and describe the evidence in
Figure~\ref{fig:results}. Explain trends, compare conditions, and connect the
results with the principles introduced earlier. Distinguish measured evidence
from proposed explanations. Replace the placeholder image with an actual
measurement, photograph, or characterization plot.

\begin{figure}[htbp]
    \centering
    \includegraphics[width=0.65\linewidth]{figC.jpg}
    \caption{Replace with a description of the observed or characterized results.}
    \label{fig:results}
\end{figure}

\subsection{Quantitative analysis and uncertainty}
Present the relevant numerical results in Table~\ref{tab:results}, with units
in the headings. Replace every bracketed entry with a measurement or a calculated
value. State the uncertainty or repeatability and explain the calculation method.

\begin{table}[htbp]
    \centering
    \caption{Replace with a description of the measured or calculated data.}
    \label{tab:results}
    \begin{tabularx}{\linewidth}{@{}lXXX@{}}
        \toprule
        Sample & Input [unit] & Response [unit] & Uncertainty [unit] \\
        \midrule
        {[Sample 1]} & {[Value]} & {[Value]} & {[Value]} \\
        {[Sample 2]} & {[Value]} & {[Value]} & {[Value]} \\
        {[Sample 3]} & {[Value]} & {[Value]} & {[Value]} \\
        \bottomrule
    \end{tabularx}
\end{table}

Use Figure~\ref{fig:comparison} to compare measurements, a fitted model, or a
reference result. Identify the axes, units, legend, and error bars in the
replacement figure. Discuss possible sources of error and limitations, and
explain whether the evidence supports the experimental objective.

\begin{figure}[htbp]
    \centering
    \includegraphics[width=0.65\linewidth]{figD.jpg}
    \caption{Replace with a description of the quantitative comparison or analysis.}
    \label{fig:comparison}
\end{figure}

\section{Conclusion}
Summarize the main findings and answer the objective stated in the Introduction.
Include the key quantitative results where appropriate. State the main
limitations and a practical improvement. Do not introduce new measurements in
this section.

% Keep one numbered References heading; the bibliography's automatic heading
% is suppressed locally to avoid the duplicate heading in the original report.
\section{References}
\begingroup
\renewcommand{\section}[2]{}%
\begin{thebibliography}{9}
    \bibitem{reference-placeholder}
    [Author(s)]. [Article or book title]. \emph{[Journal or publisher]}.
    [Year]; [Volume]([Issue]): [Pages]. [DOI or URL, if available].
\end{thebibliography}
\endgroup

\section{Questions}
\begin{enumerate}
    \item \textbf{[Insert the assigned question.]}

    \textbf{Answer:} Replace with a reasoned answer. Include a derivation,
    evidence, or a literature citation as appropriate.
\end{enumerate}

\end{document}
